\chapter{From Classical Models to Neural Networks}

\section{Limitations of Shallow Models}

In previous chapters, we studied a range of classical models, including linear regression, logistic regression, and kernel methods. These models are powerful and well-understood, yet they exhibit fundamental limitations.

\medskip

\noindent
\textbf{Shallow models.}  
We refer to models of the form
\[
f(x) = \langle w, \phi(x) \rangle,
\]
where $\phi$ is either fixed or involves only a single layer of transformation.

\medskip

\noindent
\textbf{Key limitation: representation.}  
The success of these models depends critically on the choice of features $\phi(x)$.

\medskip

\noindent
\textbf{Challenges.}
\begin{itemize}
    \item Designing effective features requires domain expertise
    \item Fixed feature maps may fail to capture complex structure
    \item High-dimensional feature spaces can be computationally expensive
\end{itemize}

\medskip

\noindent
\textbf{Expressivity limitations.}
\begin{itemize}
    \item Linear models produce linear decision boundaries
    \item Kernel methods provide flexibility, but scale poorly with data size
\end{itemize}

\medskip

\noindent
\textbf{Curse of dimensionality.}  
Explicit feature expansions often require exponentially many features to capture complex patterns.

\medskip

\noindent
\textbf{Insight.}  
Classical models separate representation and learning: features are designed first, and then a model is fitted. This separation can limit their effectiveness.

\section{Compositional Structure of Functions}

Many real-world functions exhibit a \emph{compositional structure}.

\medskip

\noindent
\textbf{Definition (informal).}  
A function is compositional if it can be written as a composition of simpler functions:
\[
f(x) = f_L(f_{L-1}(\cdots f_1(x))).
\]

\medskip

\noindent
\textbf{Examples.}
\begin{itemize}
    \item Images: edges $\rightarrow$ shapes $\rightarrow$ objects
    \item Language: characters $\rightarrow$ words $\rightarrow$ sentences
\end{itemize}

\medskip

\noindent
\textbf{Hierarchical structure.}  
Each layer captures patterns at a different level of abstraction.

\medskip

\noindent
\textbf{Efficiency of composition.}  
Compositional representations can describe complex functions using relatively few parameters.

\medskip

\noindent
\textbf{Contrast with shallow models.}  
A shallow model may require exponentially many units to represent a function that can be expressed compactly using multiple layers.

\medskip

\noindent
\textbf{Insight.}  
Depth allows reuse of intermediate computations, leading to more efficient representations.

\section{Motivation for Deep Learning}

Deep learning arises from the idea of learning both representations and functions through composition.

\medskip

\noindent
\textbf{Neural networks.}  
A neural network with $L$ layers has the form
\[
f(x) = f_L(f_{L-1}(\cdots f_1(x))),
\]
where each layer is typically an affine transformation followed by a nonlinearity.

\medskip

\noindent
\textbf{Key advantages.}

\begin{itemize}
    \item \textbf{Automatic feature learning:}  
    Features are learned from data rather than manually designed.

    \item \textbf{Expressivity:}  
    Deep networks can represent highly complex functions.

    \item \textbf{Scalability:}  
    Efficient training with large datasets and modern hardware.
\end{itemize}

\medskip

\noindent
\textbf{Universal approximation (informal).}  
Neural networks can approximate a wide class of functions, given sufficient capacity.

\medskip

\noindent
\textbf{Practical success.}  
Deep learning has achieved strong performance in:
\begin{itemize}
    \item computer vision,
    \item natural language processing,
    \item speech recognition,
    \item generative modeling.
\end{itemize}

\medskip

\noindent
\textbf{Key shift.}  
Instead of designing features, we design \emph{architectures} that can learn features automatically.

\section{Transition to Learned Representations}

The transition from classical models to deep learning involves a change in how representations are constructed.

\medskip

\noindent
\textbf{Classical approach.}
\begin{itemize}
    \item Feature map $\phi$ is fixed
    \item Model learns parameters $w$
\end{itemize}

\medskip

\noindent
\textbf{Deep learning approach.}
\begin{itemize}
    \item Feature map $\phi_\theta$ is parameterized
    \item Representation and model are learned jointly
\end{itemize}

\medskip

\noindent
\textbf{Model structure.}
\[
f(x) = g_\theta(\phi_\theta(x)),
\]
where both $\phi_\theta$ and $g_\theta$ depend on parameters.

\medskip

\noindent
\textbf{Hierarchical representations.}
\begin{itemize}
    \item Early layers: local or low-level features
    \item Intermediate layers: patterns and structures
    \item Final layers: task-specific representations
\end{itemize}

\medskip

\noindent
\textbf{Geometric viewpoint.}  
Deep networks transform data through successive nonlinear mappings, reshaping its geometry to make learning easier.

\medskip

\noindent
\textbf{Optimization perspective.}  
Learning involves optimizing a highly nonconvex objective over all layers simultaneously.

\medskip

\noindent
\textbf{Regularization through structure.}  
Architectural choices impose inductive biases that guide learning:
\begin{itemize}
    \item convolutional layers encode locality,
    \item attention mechanisms encode relationships,
    \item residual connections stabilize training.
\end{itemize}

\medskip

\noindent
\textbf{Insight.}  
Deep learning unifies representation, modeling, and optimization into a single framework.

\section{A Unifying Perspective}

We can now reinterpret the evolution of models:

\begin{itemize}
    \item Classical models: fixed representations + simple models
    \item Kernel methods: implicit high-dimensional representations
    \item Deep learning: learned hierarchical representations
\end{itemize}

\medskip

\noindent
\textbf{Core idea.}  
The power of modern machine learning lies in learning representations that make the problem easier.

\medskip

\noindent
\textbf{Key components.}
\begin{itemize}
    \item Representation (features)
    \item Function (mapping)
    \item Optimization (learning)
\end{itemize}

\medskip

\noindent
Deep learning integrates these components in a unified and scalable way.

\section{Outlook}

This chapter motivated the transition from classical models to neural networks:
\begin{itemize}
    \item limitations of shallow models,
    \item importance of compositional structure,
    \item advantages of deep architectures,
    \item shift toward learned representations.
\end{itemize}

\medskip

In the next chapter, we begin our study of deep learning by examining neural networks as compositional function approximators.

\medskip

\noindent
\textbf{Guiding principle.}  
Complex problems often require structured representations, and depth provides a natural way to build them.